Un dipôle est formé d'une charge $+q$ et d'une charge $-q$, avec $q = \qty{\v{q}}{nC}$, séparées de $2a = \qty{\fpeval{2*\v{a}}}{cm}$. Le point $P$ est sur la médiatrice du dipôle, à \qty{\v{y}}{cm} du milieu.

\begin{center}
\begin{tikzpicture}
    \draw[dashed] (0,-0.5) -- (0,3.3);
    \fill (-1.5,0) circle (3pt) node[below=3pt] {$+q$};
    \fill (1.5,0) circle (3pt) node[below=3pt] {$-q$};
    \draw (-1.5,0) -- (1.5,0);
    \draw (0,3) node {$\times$} node[right=3pt] {$P$};
    \draw[|-|] (-1.5,-0.7) -- node[below] {$2a$} (1.5,-0.7);
    \draw[|-|] (-0.3,0) -- node[left] {$y$} (-0.3,3);
\end{tikzpicture}
\end{center}

\begin{enumerate}[label=\alph*)]
    \item Déterminez le module du champ électrique produit en $P$ par chacune des charges.
    \item Déterminez le module et l'orientation du champ électrique résultant en $P$.
    \item Quelle serait la force (module et sens) sur une charge de \qty{\v{q0}}{nC} placée en $P$ ?
\end{enumerate}
